\section{The published theorems the proof rests on}\label{sec:literature}

The archive does not assume the literature; it formalises it. One file collects
\literaturerows{} theorems, each returning a named statement, and each is proved inside the
archive. All of them are used by the final theorem (\cref{sec:offpath}). \Cref{app:rows} lists
them, with what each Lean statement asserts and its sources. This section describes the
collection.

\subsection{What the file contains}
Two statements are entered twice, so there are \rowsdistinct{} distinct statements. One of
them is the conclusion itself, non-vanishing of every Dirichlet $L$-function to the right of
$\tfrac12$, returned under the hypothesis of \cref{st:vwin}. One is an empty slot, the
proposition \lean{True}. Three are statements in the archive's own terms that correspond to no
single published theorem. The remaining \rowspublished{} are formal statements of published
results. We compared each with the theorem its name refers to. The comparison rests on the Lean
statement, which we read, and on the published theorem as we know it; we did not open the
papers for it. \Cref{app:rows} gives the reason for each verdict, so that each can be checked:
\begin{center}
\begin{tabular}{@{}lr@{}}
\toprule
The published theorem, in the generality of a cited source & \rowsgeneral\\
A special case of it, or a weaker form, restricted to what the proof uses & \rowsspecial\\
Not a statement of the cited literature as it stands & \rowsother\\
\bottomrule
\end{tabular}
\end{center}
The restricted statements follow a pattern. Results on automorphic forms are stated for totally
complex base fields or for unitary groups of CM extensions; results on real groups for
$\U(p,q)$; results on $\GL_n$ over local fields in characteristic zero; theorems about
reductive groups are stated for $\GL_n$, and depth-zero cases replace general ones; results on
the characters of $\GL_n(\F_q)$ are stated for prime fields. The Galois representations
attached to automorphic representations are stated with only part of the compatibility away
from $\ell$ that the literature proves: equal dimension, equality when the local parameter is
a sum of characters with trivial monodromy, and agreement on inertia for tame, square-free
parameters at large primes. The three statements that differ are a
global consequence of the local theorem of Gelfand and Kazhdan, a form of the rapid decay of
cusp forms in Sobolev norms, and a statement named after the local zeta integrals of the
doubling method that contains the global identity instead. The proof was evidently formalised
to the extent it needs and no further. A reader who wants to use one of these formalisations
elsewhere should read its statement first.

\subsection{By area}
\emph{Class field theory and Galois cohomology.} Global Artin reciprocity
\cite{takagi1920theorie,artin1927beweis,chevalley1940theorie,tate1967global}, ray class fields
and characters, the Hasse norm theorem \cite{hasse1924darstellbarkeit,hasse1931beweis}, class
field theory for quadratic extensions \cite{hilbert1899theorie}, Chebotarev's theorem
\cite{chebotarev1926bestimmung}, a lemma of Chevalley on units \cite{chevalley1951deux}; local
and global Tate duality, the Euler characteristic formulas and Poitou--Tate duality
\cite{tate1963duality,poitou1966remarques,milne1986arithmetic}, and the formula of Greenberg
and Wiles \cite{wiles1995modular,darmon1995fermat}. Mathlib does not contain global class field
theory; the archive proves it.

\emph{Complex multiplication and Hecke characters.} The theory of the $j$-invariant and the
main theorem of complex multiplication for imaginary quadratic fields
\cite{hasse1927neue,shimura1971introduction}; Hecke's $L$-functions and Tate's thesis
\cite{hecke1918neue,hecke1920neue,tate1967fourier}; the $\ell$-adic characters of algebraic
Hecke characters and their $p$-adic Hodge theory \cite{serre1968abelian}; the rationality
theorem of Siegel and Klingen \cite{klingen1962werte,siegel1969berechnung} and the integrality
theorem of Deligne and Ribet \cite{deligne1980values,cassounogues1979valeurs}; curves with
complex multiplication and minimal discriminant \cite{gross1980arithmetic,gross1982minimal}.

\emph{Representations of $p$-adic and finite groups.} The theory of Bernstein and Zelevinsky
\cite{bernstein1976representations,bernstein1977induced,zelevinsky1980induced}, Bernstein's
irreducibility theorem \cite{bernstein1984pinvariant}, the generic unitary dual
\cite{tadic1986classification}, the Langlands classification \cite{silberger1978langlands},
Iwahori-fixed vectors \cite{borel1976admissible,casselman1980unramified}, depth zero
\cite{moy1994unrefined,moy1996jacquet}, the local converse theorem
\cite{henniart1993caracterisation}, the local Langlands correspondence for $\GL_n$ as a list of
laws \cite{harris2001geometry,henniart2000preuve,scholze2013local}, the splitting of the
metaplectic cover for unitary dual pairs \cite{kudla1994splitting}, and Green's cuspidal
characters of $\GL_n(\F_p)$ \cite{green1955characters,gelfand1970representations}.

\emph{Representations of real groups.} Harish-Chandra's discrete series for $\U(p,q)$ with
their lowest $K$-types and formal degrees
\cite{harishchandra1965discrete,harishchandra1966discrete,schmid1975some}, the theorem of
Casselman and Wallach \cite{casselman1989canonical,wallach1983asymptotic}, the factorisation
theorem of Dixmier and Malliavin \cite{dixmier1978factorisations}, Schur orthogonality for
square-integrable representations \cite{godement1947orthogonalite}, Howe duality
\cite{howe1989remarks,howe1989transcending}, Weyl's branching rule \cite{weyl1946classical},
and the metaplectic representation \cite{weil1964certains,shale1962linear}.

\emph{Automorphic forms and Galois representations.} Finiteness and rapid decay of cusp forms
\cite{harishchandra1968automorphic,gelfand1969representation,moeglin1994decomposition},
reduction theory and strong approximation
\cite{borel1962arithmetic,kneser1965starke,platonov1969problem}, the tensor product theorem
\cite{flath1979decomposition}, multiplicity one and the theorem of Jacquet and Shalika
\cite{shalika1974multiplicity,piatetskishapiro1979multiplicity,jacquet1981eulerI,jacquet1981eulerII},
purity and cohomologicality \cite{clozel1990motifs}, base change from unitary groups
\cite{labesse2011changement}, and the Galois representations attached to polarised regular
algebraic cuspidal representations
\cite{shin2011galois,chenevier2013construction,caraiani2012local}.

\emph{The theta correspondence.} Weil's convergence criterion and the Siegel--Weil formula for
unitary groups \cite{weil1965formule,ichino2007siegel,yamana2011siegel,ikeda1996residue}, the
mass formula \cite{siegel1935analytische,braun1941theorie}, the doubling method
\cite{piatetskishapiro1987lfunctions,gelbart1987explicit}, the continuation of Eisenstein
series \cite{langlands1976functional,bernstein2024meromorphic}, and the Rallis inner product
formula \cite{rallis1984injectivity,li1992nonvanishing}.

\emph{$p$-adic Hodge theory.} Fontaine's rings
\cite{fontaine1982certains,fontaine1994corps,fontaine1994semistables}, the $p$-adic monodromy
theorem \cite{berger2002representations,andre2002filtrations,mebkhout2002analogue,kedlaya2004padic}, the local
theory of Bloch and Kato \cite{bloch1990lfunctions}, prismatic $F$-crystals and the prismatic
Dieudonn\'e functor \cite{bhatt2023prismatic,anschutz2023prismatic}, the deformation theory of
Grothendieck and Messing \cite{messing1972crystals,mazur1974universal}, and Fargues's theorem on
the Hodge--Tate map \cite{fargues2008isomorphisme,tate1967pdivisible}.

\emph{Algebraic and arithmetic geometry.} Resolution of singularities and principalisation
\cite{hironaka1964resolution,kollar2007lectures}, the vanishing theorems of Kawamata--Viehweg
and Nadel \cite{kawamata1982generalization,viehweg1982vanishing,nadel1990multiplier}, the
separation of points \cite{angehrn1995effective}, extension with $L^2$ estimates
\cite{ohsawa1987extension,demailly2016extension}, the numerical theory of nef bundles
\cite{kleiman1966toward,lazarsfeld2004positivity}, the comparison of algebraic and analytic
sections \cite{serre1956geometrie}, formal algebraisation \cite{grothendieck1961elements},
abelian schemes and their duals \cite{mumford1965geometric,mumford1970abelian}, semistable
reduction \cite{grothendieck1972groupes}, toroidal compactifications
\cite{ash1975smooth,faltings1990degeneration,lan2013arithmetic,pink1989arithmetical},
arithmetic intersection theory, the arithmetic Hilbert--Samuel theorem and arithmetic volumes
\cite{gillet1990arithmetic,gillet1992arithmetic,bost1994heights,abbes1995theoreme,zhang1995positive,moriwaki2009continuity},
Chow forms \cite{chow1937algebraischen}, the Poincar\'e--Lelong formula
\cite{lelong1957integration}, semi-negativity of Lie algebras \cite{bost2004germs}, good
metrics \cite{mumford1977hirzebruch}, and Zarhin's trick \cite{zarhin1977endomorphisms}.

\emph{Other.} Linnik's theorem on the least prime in a progression
\cite{linnik1944least1,linnik1944least2}, an estimate for $L'/L(1,\chi)$
\cite{davenport1980multiplicative}, Brauer's induction theorem \cite{brauer1947artin},
Radford's theorem on shuffle algebras \cite{radford1979natural}, a simplicity theorem for
induced modules in characteristic $p$ \cite{andersen1980strong,jantzen2003representations},
the Hasse principle for Hermitian forms \cite{landherr1936aquivalenz,jacobowitz1962hermitian},
and the theorem of Stone and von Neumann \cite{stone1930linear,vonneumann1931eindeutigkeit}.

\subsection{What this means for the result}
The certificate for the Riemann Hypothesis does not depend on any of these names being
accurate: the kernel checked a closed term of type \lean{RiemannHypothesis}, whatever its
intermediate statements are called. But the archive's worth beyond that one theorem depends on
them a great deal. If the formal statements are what \cref{app:rows} reports, the archive
contains verified proofs of \rowspublished{} published theorems, \rowsgeneral{} of them in the
generality of their sources. That is a second result, separable from the first, and each row
can be checked independently against the published theorems it names.
